\documentclass[10pt,letterpaper,sans]{moderncv}
\moderncvstyle{banking}
\moderncvcolor{blue}

\usepackage[scale=0.81]{geometry}
\nopagenumbers{}
\setlength{\hintscolumnwidth}{2.25cm}

\name{Krishnaa}{Vadivel}
\title{Scientific Computing | Machine Learning | Electrical Engineering}
\email{srikrishnaa.careers@gmail.com}
\phone[mobile]{516-853-5663}
\social[linkedin]{sri-krishnaa-ece-phd}
\extrainfo{\href{https://krishnaa423.github.io/personal-website/}{Website} \textbar{} \href{https://github.com/krishnaa423}{GitHub: krishnaa423} \textbar{} \href{https://leetcode.com/u/krishnaa42342}{LeetCode: krishnaa42342} \textbar{} \href{https://hub.docker.com/u/krishnaa42342}{Docker Hub: krishnaa42342}}

\begin{document}
\makecvtitle

\section{Summary}
\cvitem{}{Electrical engineering PhD researcher with experience spanning quantum and many-body theory, first-principles materials simulation, integrated photonics, scientific machine learning, embedded systems, and distributed computing. Developed theoretical models and implemented software across Python/C++, CUDA/ROCm, FPGA/MCU platforms, web-server and database systems, container tools, cloud frameworks such as gcloud, and leadership-class HPC environments.}

\section{Experience}
\cventry{2021--present}{PhD Research Assistant}{Yale University, Electrical Engineering}{}{}{\begin{itemize}%
\item Developed machine-learning and automatic-differentiation models in Python for material-property prediction and optical-emission calculations.
\item Developed first-principles and quantum many-body simulation software in Python and C++ using MPI, OpenMP, PETSc, SLEPc, CUDA, ROCm, and CMake across leadership-class HPC systems.
\item Built distributed applications for materials and molecular simulations on Frontier, Frontera, Perlmutter, and related national-lab GPU clusters.
\item Helped fabricate and characterize semiconductor materials during the early PhD, alongside simulation and modeling work tied to photonic and electronic materials.
\item Built agentic AI pipelines for reading papers and software documentation, then generating structured inputs for first-principles simulation workflows.
\item Co-authored a 2025 \textit{Journal of the American Chemical Society} paper and delivered APS presentations in 2024, 2025, and 2026 on self-trapped exciton and exciton-polaron research.
\end{itemize}}
\cventry{2019--2021}{Software and Embedded Systems Engineer}{Probe Technologies Inc. / Weatherford Wireline Services}{}{}{\begin{itemize}%
\item Built parallel applications with CUDA and OpenCL fallback for large sensor-data and waveform processing, diagnostic analytics, and field-data visualization in deployed well-logging systems.
\item Developed full-stack server applications with ASP.NET and Microsoft SQL Server for real-time control of company tool testing, data collection, and data processing.
\item Developed C\# GUI applications for preparing customer inputs and interacting with deployed hardware systems.
\item Implemented embedded and digital functionality across Lattice FPGA, PIC32, dsPIC, STM32, and C8051 platforms, including circuit bring-up, oscilloscope validation, and targeted Altium PCB updates.
\item Used \texttt{dolfinx}-based finite-element analysis for frequency analysis of steel pipes in wells.
\end{itemize}}
\cventry{2018--2019}{Photonics Technical Writer}{FindLight}{}{}{\begin{itemize}%
\item Wrote technical articles on laser bench systems, optics, and photonics-device applications for engineering-facing audiences.
\end{itemize}}

\section{Selected Projects}
\cvitem{Machine Learning Models for Material Properties}{Developed PyTorch Geometric graph neural networks and automatic-differentiation models to accelerate material-property and optical-emission calculations in molecular and condensed-matter systems.}
\cvitem{Integrated Photonics Bandpass Filter}{Developed Meep FDTD simulations for a ring-resonator-based integrated-photonics bandpass filter that was later fabricated and tested with an erbium-doped fiber laser.}
\cvitem{Container Orchestration for First-Principles Computing}{Developed CUDA-ready and cross-environment containers for first-principles software stacks and published them to Docker Hub for portable scientific computing.}
\cvitem{Agentic AI for Scientific Workflows}{Built AI pipelines for reading documentation and journal papers, then generating structured inputs and starter artifacts for first-principles and simulation software.}

\section{Education}
\cventry{2021--present}{PhD, Electrical Engineering}{Yale University}{}{}{}
\cventry{2023}{MPhil, Electrical Engineering}{Yale University}{}{}{}
\cventry{2023}{MS, Electrical Engineering}{Yale University}{}{}{}
\cventry{2017--2018}{MEng, Electrical and Computer Engineering}{Cornell University}{}{}{}
\cventry{2013--2017}{BS, Electrical and Computer Engineering}{Cornell University}{}{}{}

\section{Technical Skills}
\cvitem{Programming}{Python, C, C++, CUDA, JavaScript, Bash, PowerShell, C\#, SQL}
\cvitem{Scientific / ML}{PyTorch, PyTorch Geometric, scikit-learn, NumPy, pandas, SciPy, SymPy, matplotlib, graph neural networks, transformers, automatic differentiation}
\cvitem{Distributed Compute}{MPI, OpenMP, PETSc, SLEPc, mpi4py, petsc4py, slepc4py, ROCm, Frontier, Frontera, Perlmutter}
\cvitem{Hardware / Systems}{Verilog, Lattice FPGA, PIC32, dsPIC, STM32, C8051, Linux, Docker, Docker Compose, Kubernetes, ASP.NET, Microsoft SQL Server, gcloud}
\cvitem{Domains}{Computational physics, semiconductor physics, integrated photonics, embedded systems, scientific software}

\end{document}
